\documentclass[11pt,a4paper]{article}

\usepackage[utf8]{inputenc}
\usepackage[T1]{fontenc}
\usepackage{helvet}
\renewcommand{\familydefault}{\sfdefault}
\usepackage[margin=2.2cm, top=2cm, bottom=2cm]{geometry}
\usepackage{graphicx}
\usepackage{xcolor}
\usepackage{booktabs}
\usepackage{float}
\usepackage{caption}
\usepackage{amsmath}
\usepackage{hyperref}
\usepackage{fancyhdr}
\usepackage{tikz}
\usepackage{enumitem}
\usepackage{tabularx}

\definecolor{bigorange}{RGB}{255,130,0}
\definecolor{biggrey}{RGB}{60,60,60}

\hypersetup{colorlinks=true, linkcolor=bigorange, urlcolor=bigorange}

\pagestyle{fancy}
\fancyhf{}
\renewcommand{\headrulewidth}{0.5pt}
\renewcommand{\headrule}{\hbox to\headwidth{\color{bigorange}\leaders\hrule height \headrulewidth\hfill}}
\fancyhead[L]{\small\color{biggrey} Banco BIG $|$ Rates Research}
\fancyhead[R]{\small\color{biggrey} May 2026}
\fancyfoot[C]{\small\color{biggrey}\thepage}

\begin{document}

% COVER
\thispagestyle{empty}
\begin{tikzpicture}[remember picture, overlay]
    \fill[bigorange] (current page.north west) rectangle ([yshift=-9cm]current page.north east);
    \fill[white, opacity=0.15] ([yshift=-3cm]current page.north west) -- ([yshift=-7cm]current page.north west) -- ([yshift=-5cm]current page.north east) -- ([yshift=-1cm]current page.north east) -- cycle;
    \fill[bigorange] ([yshift=1.5cm]current page.south west) rectangle (current page.south east);
\end{tikzpicture}

\vspace*{1cm}
\begin{center}
    
\includegraphics[width=4.5cm]{../Logo-BIG.png}
\end{center}

\vspace{3.5cm}

\begin{center}
    {\Huge\bfseries\color{biggrey} Interest Rate Spillovers}\\[0.4cm]
    {\Huge\bfseries\color{biggrey} to the US Treasury Curve}\\[1.0cm]
    {\Large\color{biggrey} 5-Country VAR: US, EU, CA, BR, UK}\\[0.3cm]
    {\Large\color{biggrey} Nelson-Siegel Level, Slope, Curvature}\\[2.0cm]
    {\large\color{biggrey} Rates Research $|$ 5 May 2026}
\end{center}

\vfill
\begin{tikzpicture}[remember picture, overlay]
    \node[anchor=south, yshift=2.0cm] at (current page.south) {
        \color{white}\small\bfseries Banco BIG $|$ Institutional Research
    };
\end{tikzpicture}

\newpage

% PAGE 2: METHODOLOGY & RESULTS
\section*{Methodology: 5-Country VAR on Nelson-Siegel Factors}

We decompose five government yield curves (US, EU, CA, BR, UK) into Level, Slope, and Curvature via Nelson-Siegel ($\lambda=0.10$). A 15-variable VAR(1) is estimated on weekly data (2005--2026, $N=1112$ obs).

\smallskip
\noindent\textbf{Cholesky ordering:} EU $\rightarrow$ CA $\rightarrow$ BR $\rightarrow$ UK $\rightarrow$ US

\smallskip
\noindent Rationale: US is placed last as the ``receiver'' of interest. US is most exogenous (global reserve currency).

\subsection*{Granger Causality Tests ($\alpha=5\%$)}

\begin{table}[H]
\centering\small
\begin{tabular}{llccc}
\toprule
\textbf{Source} & \textbf{US Factor} & \textbf{F-stat} & \textbf{p-value} & \textbf{Sig.?} \\
\midrule
    EU & Level & 7.07 & 0.000 & \textbf{Sig.} \\
    EU & Slope & 7.35 & 0.000 & \textbf{Sig.} \\
    EU & Curvature & 3.41 & 0.017 & \textbf{Sig.} \\
    CA & Level & 0.85 & 0.466 & Not sig. \\
    CA & Slope & 0.33 & 0.801 & Not sig. \\
    CA & Curvature & 0.84 & 0.470 & Not sig. \\
    BR & Level & 0.54 & 0.652 & Not sig. \\
    BR & Slope & 0.97 & 0.406 & Not sig. \\
    BR & Curvature & 0.37 & 0.773 & Not sig. \\
    UK & Level & 6.70 & 0.000 & \textbf{Sig.} \\
    UK & Slope & 7.08 & 0.000 & \textbf{Sig.} \\
    UK & Curvature & 4.15 & 0.006 & \textbf{Sig.} \\
\bottomrule
\end{tabular}
\end{table}

\begin{figure}[H]
    \centering
    \includegraphics[width=\textwidth]{fig_factors.png}
    \vspace{-0.3cm}
    \caption{\small Nelson-Siegel factors for all 5 countries (2005--2026). US highlighted as the target curve.}
\end{figure}

\newpage

% PAGE 3: IRFs AND FEVD
\section*{Shock Propagation to United States}

\begin{figure}[H]
    \centering
    \includegraphics[width=\textwidth]{fig_irf.png}
    \vspace{-0.3cm}
    \caption{\small Orthogonalised IRFs: US factor responses to foreign shocks over 40 weeks. Blue shading = 95\% bootstrap CI.}
\end{figure}

\subsection*{FEVD at 52-Week Horizon (\% of variance)}

\begin{table}[H]
\centering\small
\begin{tabular}{l c c c c c}
\toprule
 & \textbf{US} & \textbf{EU} & \textbf{CA} & \textbf{BR} & \textbf{UK} \\
\midrule
    US Level & 36.1 & 14.3 & 28.3 & 5.0 & 16.4 \\
    US Slope & 35.6 & 14.5 & 22.4 & 0.9 & 26.6 \\
    US Curvature & 26.9 & 23.6 & 33.2 & 5.2 & 11.0 \\
\bottomrule
\end{tabular}
\end{table}

\begin{figure}[H]
    \centering
    \includegraphics[width=\textwidth]{fig_fevd.png}
    \vspace{-0.3cm}
    \caption{\small FEVD: share of US factor forecast error variance by source country at horizons 1--52 weeks.}
\end{figure}

\newpage

% PAGE 4: MARKET CONTEXT
\section*{United States: Market Context (May 2026)}


\subsection*{Macro \& Fiscal Backdrop}
\begin{itemize}[nosep, leftmargin=1.2em]
    \item \textbf{Inflation:} Core PCE at 2.6\% (Mar 2026). Energy pass-through from Iran conflict adding $\sim$0.5pp to headline.
    \item \textbf{Growth:} Q1 2026 GDP at 1.8\% annualised, showing resilience despite tighter financial conditions.
    \item \textbf{Fed:} Held at 4.25--4.50\%. Markets no longer pricing cuts in 2026; $\sim$2 hikes now priced.
    \item \textbf{Issuance:} US Treasury running $\sim$\$2tn annual deficit financing. Refunding mix shifting toward bills.
    \item \textbf{Debt ceiling:} Ongoing uncertainty adds term premium pressure.
\end{itemize}

\subsection*{Key Drivers Ahead}
\begin{enumerate}[nosep, leftmargin=1.4em]
    \item \textbf{US is the global anchor:} As the most exogenous rate in our VAR, US shocks propagate to all other curves. Own-variance share ($\sim$36\%) is relatively low, reflecting deep global integration.
    \item \textbf{Canada co-movement:} CA explains 28\% of US Level variance---reflecting NAFTA integration and BoC/Fed policy synchronisation.
    \item \textbf{Fiscal trajectory:} CBO projections of rising debt/GDP create persistent term premium.
    \item \textbf{Iran ceasefire durability:} Key swing factor for energy pass-through and Fed reaction function.
\end{enumerate}


\vspace{0.3cm}
\noindent\rule{\textwidth}{0.4pt}
{\footnotesize\color{biggrey} \textbf{Data:} US: FRED (constant-maturity Treasury rates), EU: ECB Statistical Data Warehouse (AAA-rated), CA: Bank of Canada Valet API (benchmark bonds), BR: IPEA/ANBIMA (DI pre-fixed yield curve), UK: Bank of England IADB (nominal spot curve). Model: VAR(1) on weekly NS factors, sample 2005--2026 (1112 obs). This document is for informational purposes only and does not constitute investment advice.}

\end{document}
